\section{The Rule}
\label{sec:rule}

Each beat applies the coin, then the meetings, then the collision, then the stream. Every piece is an exact,
reversible map in $\mathbb{Z}[\omega][1/2]$ with no float, no rounding and no random number, and runs backward piece
by piece. Each piece below is the one that survived a stated alternative.

\begin{itemize}
\item \textbf{The knit} is the whole set of pieces, and it keeps $W(F_4)$, the full symmetry of the dock (1,152
elements). Such a knit exists and keeps CPT, so no direction of the dock is preferred (\code{E-RLT-0061}).
\item \textbf{The collision.} A dock's full lines turn as $-1$ and a lone vibe goes straight through, so a lone
vibe's wake stays on its own line, the free particle's straight path (\code{E-RLT-0084}).
\item \textbf{The doublet lock.} A vibe's take direction along its line is its role's spin doublet. It is the only
way a vibe can move with amplitudes and stay covariant, and it makes every charge-one state a spinor
(\code{E-RLT-0097}, \code{E-SPN-0071}).
\item \textbf{The meeting.} A love and a fear on one line pass unchanged. Two like vibes keep their points (weight
$1/4$) or exchange them ($3/4$), leaving a \emph{knot} with CHSH $\sqrt7$ (\code{E-RLT-0099}).
\item \textbf{The contact.} Two like vibes on a full line keep their slots, the \emph{pass}, phase $+1$, instead of
bouncing, $-1$. Of the six phases the ring allows, only these two add no new number, and only the pass binds the
electron (\code{E-SPN-0092}, \code{E-SPN-0093}).
\item \textbf{The store.} Pairs are made and released with no veto, each stored pair holding two points. It is the
only vacuum measured that holds under the pass (\code{E-RLT-0102}, \code{E-RLT-0103}).
\item \textbf{The coin.} On a line holding one open vibe, $2C = (1 + \omega) I + (1 - \omega) X$ on its two slots.
It gives a lone vibe a mass, with a gap of $2\pi/3$ and a top speed of $1/2$ a beat, and its cost is exact lattice
momentum, the one symmetry that forbids a mass (\code{E-SPN-0090}, \code{E-SPN-0091}).
\end{itemize}

\subsection{The working vacuum}

The vacuum is the no-veto store under the pass, with the coin on, and one C domain with no walls. On 51 of 51 gated
cases (\code{E-RLT-0105}) it keeps its conservation laws, because no vacuum vibe is ever alone on its line and the
coin acts on it 0 times on every path. It runs backward exactly, keeps C, and makes every husk dock one causal world.
Pairs made equal pairs unmade, 466,944 each, on a cycle of 12 beats (\code{E-GRV-0072}). A lone love spreads over 105
to 124 occupations by beat 16 on 17 of 17 starts, so positions carry amplitude at no cost to the vacuum. Not held:
the kept branch of a like meeting does not interfere, and the lone love reads the dephased $1/4$, $1/16$, $1/64$.

\subsection{The line law}

With no mixer, the tone on every straight mesh line is conserved: 6,144 of 6,144 lines on side 8 over 48 beats
(\code{E-SPN-0098}). Every piece acts inside one line, so the knit is 6,144 one-dimensional worlds that never touch.
Line locality is momentum conservation: a lone vibe's momentum is its root, so it cannot change line without changing
momentum. The vacuum is therefore correlated without limit along each line and not at all across lines, and three
failures are one fact: no $1/r$ response (\code{E-GRV-0086}), no area law (\code{E-GRV-0083}), and no quiet reach
across lines (\code{E-SPN-0101}). Every mixer the rule's ring allows cascades through this vacuum
(\code{E-SPN-0094} to \code{E-SPN-0099}), and no store can join lines, since an involution must return every vibe to
its slot (\code{E-RLT-0107}). Section~\ref{sec:motion} is where this is resolved.
